% ----------------------------------------------------------------------------
\chapter{Flows, transversality and the total curvature bound}\label{ch:flows}
\module{NoCompromise.Flow.FlowBox,
        NoCompromise.Flow.ODE,
        NoCompromise.Surface.Geometry,
        NoCompromise.Surface.Morse,
        NoCompromise.Surface.MorseCount,
        NoCompromise.Surface.Shape,
        NoCompromise.Surface.TotalCurvature,
        NoCompromise.Topology.Parity,
        NoCompromise.Topology.Transversality,
        NoCompromise.Topology.Tubular}

This chapter supplies the differential-topological layer.  Its output is used
in exactly three places: Theorem~\ref{thm:component-count} by
Lemma~\ref{lem:hull-properties} and Lemma~\ref{lem:level-connected},
Theorem~\ref{thm:tubular} by Lemma~\ref{lem:interior-flux}, and the total
curvature bound \eqref{eq:total-curvature-bound} by
Lemma~\ref{lem:K-gauss-bonnet-input}.

It is the most expensive part of the project
(Formalisation note~\ref{fnote:flow-cost}).  Formalisation
note~\ref{fnote:why-degree} records what the present design already avoids,
and Formalisation note~\ref{fnote:retire-flows} what might still be avoided.

\section{From Picard--Lindel\"of to a global smooth flow}\label{sec:flow}

\begin{theorem}[Global flow]\label{thm:global-flow}
\uses{conv:ambient}
Let $X:\R^n\to\R^n$ be $C^k$, $1\le k\le\infty$, globally Lipschitz with
constant $L$ and bounded by $M$.  Then there is a unique
$\Phi:\R\times\R^n\to\R^n$ with
\begin{equation}\label{eq:flow-integral}
  \Phi_t(x)=x+\int_0^tX(\Phi_s(x))\,ds .
\end{equation}
In particular a compactly supported $C^k$ field satisfies these hypotheses,
with $M=\|X\|_\infty$ and $L=\|DX\|_\infty$.
\end{theorem}

\begin{proof}
The local Picard--Lindel\"of family \ref{base:ode} applies on every closed
ball: continuity supplies $M$, the mean value theorem supplies the spatial
Lipschitz hypothesis with constant $L$.  Its remaining hypothesis is the
interval condition
$M\cdot\max(t_{\max}-t_0,\,t_0-t_{\min})\le a-r$ relating the time interval to
the ball radii $r<a$ on which the family is posed.  It holds on an
arbitrarily long interval: given $T>0$, take $r=a/2$ with $a\ge2MT$, so that
$a-r=a/2\ge MT$.  \emph{This is the only place global boundedness is used, and
it is what upgrades a local statement to a flow defined for all time.}

A local solution obeys $|\Phi_t(x)-x|\le M|t|$, so on a prescribed bounded
time interval its graph stays in a fixed compact ball.  The quantitative local
theorem supplies a time step depending only on the bounds of $X$ on that ball;
finitely many consecutive steps cover the interval, and uniqueness makes
consecutive solutions agree and extensions from different intervals
compatible.
\end{proof}

\begin{lemma}[Group law]\label{lem:flow-group}
\uses{thm:global-flow}
$\Phi_0=\Id$, $\Phi_{t+s}=\Phi_t\circ\Phi_s$, $\Phi_t^{-1}=\Phi_{-t}$.
\end{lemma}

\begin{proof}\uses{thm:global-flow}
Uniqueness in \eqref{eq:flow-integral} after a time translation.
\end{proof}

\begin{lemma}[Gr\"onwall and bi-Lipschitz dependence]\label{lem:flow-gronwall}
\uses{thm:global-flow}
$|\Phi_t(x)-\Phi_t(y)|\le e^{L|t|}|x-y|$; hence each $\Phi_t$ is a
bi-Lipschitz homeomorphism.
\end{lemma}

\begin{proof}\uses{thm:global-flow,lem:flow-group}
Subtract \eqref{eq:flow-integral} at $x$ and $y$ and apply the elementary
Gr\"onwall inequality, proved by differentiating
$e^{-Lt}\int_0^tL|\Phi_s(x)-\Phi_s(y)|\,ds$, first for $t\ge0$ and then
backwards in time.
\end{proof}

\begin{lemma}[The variational equation]\label{lem:flow-variational}
\uses{thm:global-flow}
For fixed $x$ the linear matrix integral equation
\begin{equation}\label{eq:flow-variational}
  A(t,x)=I+\int_0^tDX(\Phi_s(x))A(s,x)\,ds
\end{equation}
has a unique solution, and $(t,x)\mapsto A(t,x)$ is continuous.
\end{lemma}

\begin{proof}\uses{thm:global-flow}
Picard iteration in the finite-dimensional matrix space converges uniformly on
compact time intervals: the difference of the $m$-th and $(m+1)$-st iterates is
bounded by $C(L|t|)^m/m!$.  The same estimate and continuity of
$DX\circ\Phi$ give joint continuity.
\end{proof}

\begin{theorem}[$C^k$ dependence on the initial point]\label{thm:flow-Ck}
\uses{thm:global-flow,lem:flow-variational}
$D_x\Phi_t(x)=A(t,x)$, continuously in $(t,x)$; each $\Phi_t$ is a $C^1$
diffeomorphism with $(D\Phi_t)^{-1}=D\Phi_{-t}(\Phi_t(x))$; and if
$X\in C^k$ then $(t,x)\mapsto\Phi_t(x)$ is $C^k$.  In particular a smooth
vector field has a smooth flow.
\end{theorem}

\begin{proof}\uses{lem:flow-variational,lem:flow-gronwall,lem:flow-group}
For an increment $z$, subtract \eqref{eq:flow-integral} at $x+z$ and $x$, then
subtract \eqref{eq:flow-variational} applied to $z$.  By
Lemma~\ref{lem:flow-gronwall} all points lie in one compact set and
$|X(p+w)-X(p)-DX(p)w|\le\omega(|w|)|w|$ with $\omega(r)\downarrow0$ a modulus
of continuity for $DX$ there.  A second Gr\"onwall estimate gives, uniformly
for $t$ in a fixed compact interval,
\begin{equation}\label{eq:flow-differentiable}
  \frac{|\Phi_t(x+z)-\Phi_t(x)-A(t,x)z|}{|z|}
  \le C_t\,\omega\bigl(e^{L|t|}|z|\bigr)\longrightarrow0 .
\end{equation}
Invertibility follows from Lemma~\ref{lem:flow-group} and the chain rule.  For
$X\in C^k$, differentiate \eqref{eq:flow-variational} inductively: at order
$j$ the highest derivative of $\Phi$ occurs linearly with coefficient
$DX(\Phi)$, while the forcing contains only derivatives of $X$ of order
$\le j$ and already-constructed lower-order derivatives of $\Phi$.  The same
integral-equation and Gr\"onwall argument gives $C^j$ dependence; induct to
$j=k$.  The ODE $\partial_t\Phi=X\circ\Phi$ supplies the time and mixed
derivatives.
\end{proof}

\begin{theorem}[Flow box]\label{thm:flow-box}
\uses{thm:flow-Ck}
Let $X\in C^k$ with $X(p)\ne0$.  Choose an affine $(n-1)$-disk $S$ through $p$
whose tangent space is complementary to $\R X(p)$, and set
\begin{equation}\label{eq:flow-box}
  \Psi:S\times(-\varepsilon,\varepsilon)\to\R^n,\qquad\Psi(y,t):=\Phi_t(y).
\end{equation}
After shrinking the domain, $\Psi$ is a $C^k$ coordinate chart, and in these
coordinates $X$ becomes $\partial_t$.
\end{theorem}

\begin{proof}\uses{thm:flow-Ck}
At $(p,0)$, $D\Psi(v,s)=v+sX(p)$, an isomorphism; apply the inverse function
theorem \ref{base:calculus}, and differentiate in $t$.
\end{proof}

\begin{corollary}[Flows on a manifold]\label{cor:flow-manifold}
\uses{thm:flow-Ck,thm:flow-box}
Let $M$ be a smooth manifold without boundary, let $U\subset M$ be open, and
let $X$ be a $C^k$ vector field on $U$, $1\le k\le\infty$.  Through every
$p\in U$ there is a unique maximal integral curve, the resulting local flow
is $C^k$ and satisfies the group law wherever both sides are defined, and the
flow-box conclusion holds wherever $X\ne0$.  A maximal curve cannot have a
finite endpoint while remaining in a compact subset of $U$.  Consequently a
$C^k$ vector field on a compact manifold is complete.
\end{corollary}

\begin{proof}\uses{thm:global-flow,thm:flow-Ck,thm:flow-box}
Perform the Euclidean local construction in charts.  Coordinate
representatives agree on overlaps by ODE uniqueness, so they glue; the
difference-quotient proof of Theorem~\ref{thm:flow-Ck} is invariant under
smooth coordinate changes and gives $C^k$ dependence.  The group law follows
from uniqueness, and Theorem~\ref{thm:flow-box} applies in a chart.
If a curve is confined to a compact subset of $U$, finitely many relatively
compact existence charts cover that set and furnish a common positive
continuation time after taking the minimum of their local times.  Thus a
finite maximal endpoint would admit an extension, a contradiction.  Take
$U=M$ when $M$ is compact.
\end{proof}

\begin{fnote}[Cost]\label{fnote:flow-cost}
\S\ref{sec:flow} is a mathlib-scale contribution on its own: global flow from
Picard--Lindel\"of, group law, Gr\"onwall, the variational equation solved by
Picard iteration in matrix space, $C^k$ by induction, and straightening.
Budget accordingly.  It is the single largest item in this chapter, and
Formalisation note~\ref{fnote:retire-flows}(a) records the one route that
might remove it.
\end{fnote}

\section{Relative transversality}\label{sec:transversality}

\begin{theorem}[Relative transversality for paths and disks]
\label{thm:transversality}
\uses{cor:sard-charts,thm:sard-4to3}
Let $N$ be a compact path (a $1$-manifold with boundary) or a compact disk,
let $\Sigma\subset\R^3$ be a compact smooth embedded surface, and let
$f:N\to\R^3$ be smooth, already transverse to $\Sigma$ on a neighbourhood
of $\partial N$, and with $f|_{\partial N}$ transverse to $\Sigma$ (for a path
this means $f(\partial N)\cap\Sigma=\emptyset$).  Then for every $\epsilon>0$
there is a smooth $g:N\to\R^3$ with $\|g-f\|_{C^1}<\epsilon$, $g=f$ on a
neighbourhood of $\partial N$, $g$ transverse to $\Sigma$, and $g|_{\partial N}$
transverse to $\Sigma$.
\end{theorem}

\begin{proof}\uses{cor:sard-charts,thm:sard-4to3}
Choose boundary collars $U_0\Subset U_1$ whose closures lie in the
neighbourhood where $f$ is transverse, and a smooth $\chi:N\to[0,1]$ vanishing
on $\overline{U_0}$ and positive on $N\setminus\overline{U_0}$ (obtained from a
smooth collar coordinate by a cutoff flat at the inner edge of $U_0$).  For
$v\in\R^3$ set
\begin{equation}\label{eq:transv-family}
  f_v(p):=f(p)+\chi(p)v .
\end{equation}
On $\{\chi>0\}$ the evaluation map $\mathcal F(p,v):=f_v(p)$ is a submersion,
since its $v$-derivative is $\chi(p)I$.  Hence
$Z:=\mathcal F^{-1}(\Sigma)\cap(\{\chi>0\}\times\R^3)$ is a smooth manifold of
dimension $\dim N+2$.  Let $\pi:Z\to\R^3$ be the parameter projection.
Directly from the tangent-space equation
\begin{equation}\label{eq:transv-tangent}
  D_pf_v(\xi)+\chi(p)w\in T_{f_v(p)}\Sigma
\end{equation}
one sees that $D\pi$ is onto at $(p,v)$ exactly when
$D_pf_v(T_pN)+T_{f_v(p)}\Sigma=\R^3$.  So a regular value $v$ of $\pi$ makes
$f_v$ transverse wherever $\chi>0$.

For a path $\pi$ has dimensions $3\to3$ and
Corollary~\ref{cor:sard-charts} makes its critical values null; for a disk it
has dimensions $4\to3$ and Theorem~\ref{thm:sard-4to3} does the same.  So every
ball about $0$ contains a regular $v$.  On $\overline{U_0}$ the map is
unchanged; elsewhere transversality follows from regularity of $v$.
\end{proof}

\begin{corollary}[Preimage structure]\label{cor:transv-preimage}
\uses{thm:transversality}
Let $g:N\to\R^3$ be transverse to $\Sigma$ with $g|_{\partial N}$ transverse
to $\Sigma$, as produced by Theorem~\ref{thm:transversality}.  A transverse
path has a compact zero-dimensional inverse image of $\Sigma$, hence finitely
many intersections.  A transverse disk has a compact one-dimensional inverse
image, a compact $1$-manifold with boundary whose boundary is exactly its
intersections on $\partial N$.  (Without transversality of $g|_{\partial N}$
the disk statement fails: the inclusion of the unit disk into $\{z=0\}$ is
transverse to the unit sphere centred at $(2,0,0)$, yet the preimage is the
single boundary point $(1,0)$.)
\end{corollary}

\begin{proof}\uses{thm:transversality}
The preimage theorem, i.e.\ the implicit function theorem \ref{base:calculus}
in a defining-function chart for $\Sigma$.
\end{proof}

\begin{lemma}[Even cardinality of a compact $1$-manifold boundary]
\label{lem:handshake}
The boundary of a compact one-dimensional manifold with boundary has even
cardinality.
\end{lemma}

\begin{proof}
Choose a finite cover by interval and half-interval charts and subdivide at
the finitely many chart endpoints.  The resulting finite graph has degree two
at every interior vertex and degree one exactly at the manifold boundary
points; the handshake identity $\sum_v\deg v=2\#\{\text{edges}\}$ makes the
number of degree-one vertices even.
\end{proof}

\section{Uniform tubular coordinates}\label{sec:tubular}

\begin{theorem}[Tubular neighbourhood]\label{thm:tubular}
\uses{conv:ambient}
Let $\Sigma\subset\R^3$ be a compact smooth embedded surface with a smooth
unit normal $n$.  Then
there is $\varepsilon>0$ such that
\begin{equation}\label{eq:tubular}
  \Theta:\Sigma\times(-\varepsilon,\varepsilon)\to\R^3,
  \qquad
  \Theta(p,t):=p+t\,n(p)
\end{equation}
is a diffeomorphism onto an open neighbourhood of $\Sigma$.
\end{theorem}

\begin{proof}
$D\Theta$ is invertible at every $(p,0)$, so the inverse function theorem
gives injective normal coordinates near each such point; compactness gives
finitely many neighbourhoods and a uniform local radius.  If no smaller
uniform radius made $\Theta$ globally injective, there would be
$p_j+t_jn(p_j)=q_j+s_jn(q_j)$ with $(p_j,t_j)\ne(q_j,s_j)$ and
$|t_j|+|s_j|\to0$; after taking limits $p_j,q_j\to p$, and for large $j$ both
pairs lie in the same injective normal chart at $p$, a contradiction.
\end{proof}

\section{Orientation by parity}

\begin{proposition}[Parity, orientation and component count]
\label{prop:orientation-parity}
\uses{thm:transversality,lem:handshake,thm:tubular}
Let $\Sigma\subset\R^3$ be a compact connected smooth embedded closed
surface.  Then:
\begin{enumerate}[label=(\roman*)]
\item\label{it:parity-well} the parity $I(x)\in\Z/2$ of the number of
  transverse intersections with $\Sigma$ of a path from a fixed far base point
  $z$ to $x\notin\Sigma$ is independent of the path;
\item\label{it:parity-normal} the two parity classes select compatible local
  sides, hence a continuous --- and by smoothness of $\Sigma$, smooth ---
  global unit normal $n$; in particular $\Sigma$ is orientable;
\item\label{it:parity-two} $\R^3\setminus\Sigma$ has exactly two connected
  components.
\end{enumerate}
\end{proposition}

\begin{proof}
\uses{thm:transversality,cor:transv-preimage,lem:handshake,thm:tubular}
\ref{it:parity-well}: A piecewise smooth path can be perturbed with endpoints
fixed, by Theorem~\ref{thm:transversality}, to meet $\Sigma$ transversely in
finitely many points.  Concatenating two paths gives a closed loop $\gamma$;
smooth its finitely many corners in small balls disjoint from the other
intersection points.  Choose $a\in\R^3$ and smooth
$\rho:[0,1]\to[0,1]$ zero near $0$ and one near $1$; the explicit map
\begin{equation}\label{eq:spanning-disk}
  F(re^{i\theta}):=(1-\rho(r))a+\rho(r)\gamma(e^{i\theta})
\end{equation}
is a smooth disk spanning $\gamma$, constant near the centre and agreeing with
$\gamma$ on a boundary collar.  Apply Theorem~\ref{thm:transversality} in its
disk case to perturb the interior while fixing that collar.  By
Corollary~\ref{cor:transv-preimage} the inverse image of $\Sigma$ is a compact
$1$-manifold whose boundary is exactly the intersections of $\gamma$ with
$\Sigma$, and Lemma~\ref{lem:handshake} makes that set even.

\ref{it:parity-normal}: $I$ is locally constant on the complement.  In a
defining-function chart around a point of $\Sigma$ the two local sides have
opposite parity, so both values occur; two local positive sides with the same
parity agree on their common connected collar.  Only after this gluing may
Theorem~\ref{thm:tubular} be invoked.

\ref{it:parity-two}: Let $x,y$ have the same parity and take a transverse path
between them, with an even number of intersections.  Pair consecutive
intersections; between the points of one pair the path lies on one local side
of $\Sigma$.  Since a connected smooth surface is path connected and remains
so after deleting finitely many points --- cover a connecting path by
coordinate disks and, whenever it passes through a forbidden point, replace
the corresponding diameter segment by a small semicircle in that disk --- join
the two intersection points by a path on $\Sigma$ avoiding the other
intersections.  Push it a small distance into the same globally labelled side
inside the tubular neighbourhood and join it to the old path by two small
collars.  This removes exactly two intersections and stays in the complement;
repeat.
\end{proof}

\begin{theorem}[Component count for a disjoint union of surfaces]
\label{thm:component-count}
\uses{prop:orientation-parity}
Let $\Sigma=\Sigma_1\mathbin{\dot\cup}\cdots\mathbin{\dot\cup}\Sigma_k\subset\R^3$
be a finite disjoint union of compact connected smooth embedded closed
surfaces.  Then $\R^3\setminus\Sigma$ has exactly $k+1$ connected components.
\end{theorem}

\begin{proof}\uses{prop:orientation-parity,thm:tubular}
Induct on $k$; the case $k=1$ is
Proposition~\ref{prop:orientation-parity}\ref{it:parity-two}.  Suppose the
result for $\Sigma_1,\dots,\Sigma_{k-1}$.  The connected surface $\Sigma_k$
lies in one connected component $U$ of their complement; $U$ is open and
connected in Euclidean space, hence path connected, and its tubular
neighbourhood about $\Sigma_k$ has positive distance from the other surfaces.
Apply the parity and path-surgery argument of
Proposition~\ref{prop:orientation-parity} to paths lying in $U$.  The spanning
disk used only to prove parity independence may cross the other surfaces; that
is irrelevant, since transversality is taken only with respect to $\Sigma_k$,
and the path surgery itself takes place in $U$.  Thus $U\setminus\Sigma_k$ has
exactly two components and all others are unchanged.
\end{proof}

\section{The shape operator and height functions}\label{sec:shape}

Everything in this section is pointwise linear algebra on a smooth surface.
No Morse theory, no flow, and no moving frame is used.

\begin{definition}[Shape operator]\label{def:frame-free}
\uses{conv:curvature}
For a smooth embedded surface $\Sigma\subset\R^3$ with unit normal $n$, the
\emph{shape operator} at $p$ is $A_pX:=D_Xn$ for $X\in T_p\Sigma$, the second
fundamental form is $\mathrm{II}_p(X,Y):=\langle A_pX,Y\rangle$, and
\[
  H(p):=\operatorname{tr}A_p,\qquad \kappa(p):=\det A_p ,
\]
consistently with Convention~\ref{conv:curvature}.  $A_p$ is symmetric, so
$\mathrm{II}_p$ is a symmetric bilinear form on the two-dimensional space
$T_p\Sigma$.
\end{definition}

\begin{lemma}[Hessian of a height function]\label{lem:hess-height}
\uses{def:frame-free}
Let $\Sigma\subset\R^3$ be a smooth embedded surface with unit normal $n$,
$v\in\Sph^2$, and $h_v:=\langle\cdot,v\rangle$.  Then $p$ is a critical point
of $h_v|_\Sigma$ if and only if $n(p)=\pm v$, and at such a point
\begin{equation}\label{eq:hess-height}
  \operatorname{Hess}_p(h_v|_\Sigma)
  =-\langle n(p),v\rangle\,\mathrm{II}_p
  =\mp\,\mathrm{II}_p ,
\end{equation}
where the upper displayed sign corresponds to
$\langle n(p),v\rangle=+1$.
\end{lemma}

\begin{proof}\uses{def:frame-free}
$d_ph_v(X)=\langle X,v\rangle$, which vanishes for all $X\in T_p\Sigma$ iff
$v\perp T_p\Sigma$ iff $n(p)=\pm v$.  At such a $p$, for $X,Y$ tangent,
\[
  \operatorname{Hess}_ph_v(X,Y)=\langle D_XY,v\rangle
  =\langle n(p),v\rangle\,\langle D_XY,n(p)\rangle
  =-\langle n(p),v\rangle\,\mathrm{II}_p(X,Y),
\]
where the middle step uses $v=\langle n(p),v\rangle\,n(p)$ (valid because
$v=\pm n(p)$), while the last follows by differentiating
$\langle Y,n\rangle=0$ and using Definition~\ref{def:frame-free}.
\end{proof}

\begin{lemma}[The Gauss map Jacobian is the Gauss curvature]\label{lem:gauss-jacobian}
\uses{def:frame-free,def:jacobian}
Let $\Sigma\subset\R^3$ be a smooth embedded surface with unit normal $n$.
Then $T_{n(p)}\Sph^2=T_p\Sigma$, the differential
$d_pn:T_p\Sigma\to T_{n(p)}\Sph^2$ equals $A_p$, and
\begin{equation}\label{eq:gauss-jacobian}
  \det(d_pn)=\det A_p=\kappa(p),
  \qquad
  J_2(d_pn)=|\kappa(p)| .
\end{equation}
\end{lemma}

\begin{proof}\uses{def:frame-free,def:jacobian}
$T_{n(p)}\Sph^2=n(p)^\perp=T_p\Sigma$ because $n(p)\perp T_p\Sigma$.  Then
$d_pn(X)=D_Xn=A_pX$ by Definition~\ref{def:frame-free}.  For a linear map
between two-dimensional inner-product spaces
Definition~\ref{def:jacobian} gives $J_2=|\det|$.
\end{proof}

\begin{convention}[Orientations for the Gauss map]\label{conv:gauss-orientation}
\uses{lem:gauss-jacobian,prop:orientation-parity}
Orient $\Sigma$ by the unit normal $n$ of
Proposition~\ref{prop:orientation-parity}\ref{it:parity-normal} and orient
$\Sph^2$ by its outward normal, which at the point $n(p)$ \emph{is} $n(p)$.
Under the canonical identification $T_p\Sigma=T_{n(p)}\Sph^2$ of
Lemma~\ref{lem:gauss-jacobian} these two orientations agree, so
\begin{equation}\label{eq:gauss-orientation}
  \varepsilon_n(p):=\operatorname{sign}\det\bigl(d_pn\bigr)
  =\operatorname{sign}\kappa(p)
\end{equation}
with no further sign.  \emph{This is the one place where a stray sign would
invert the conclusion of \S\ref{sec:total-curvature}}; it is recorded as a
convention so that it is fixed once.
\end{convention}

\begin{lemma}[Nondegeneracy and the sign of the curvature]
\label{lem:sign-kappa-index}
\uses{lem:hess-height,lem:gauss-jacobian}
Let $v\in\Sph^2$ and let $p$ be a critical point of $h_v|_\Sigma$.  Then $p$
is nondegenerate if and only if $\kappa(p)\ne0$, and in that case
\begin{equation}\label{eq:sign-kappa-index}
  \operatorname{sign}\kappa(p)=(-1)^{\mathrm{ind}_p(h_v)} .
\end{equation}
\end{lemma}

\begin{proof}\uses{lem:hess-height,lem:gauss-jacobian}
By \eqref{eq:hess-height},
$\det\operatorname{Hess}_ph_v=\det(\pm\mathrm{II}_p)=\det\mathrm{II}_p=\kappa(p)$,
the sign disappearing because the form is two-dimensional.  So $p$ is
nondegenerate iff $\kappa(p)\ne0$.  For a nondegenerate critical point of a
function on a surface, $\operatorname{sign}\det\operatorname{Hess}
=(-1)^{\mathrm{ind}}$: indices $0$ and $2$ give two eigenvalues of the same
sign, index $1$ two of opposite signs.
\end{proof}

\begin{lemma}[Critical points split by normal direction]\label{lem:critical-normal}
\uses{lem:hess-height,lem:gauss-jacobian}
Let $\Sigma\subset\R^3$ be closed and embedded with unit normal $n$ and let
$v\in\Sph^2$.  Then
\begin{equation}\label{eq:crit-split}
  \mathrm{Crit}(h_v|_\Sigma)=n^{-1}(v)\ \sqcup\ n^{-1}(-v),
\end{equation}
and $h_v|_\Sigma$ is a Morse function if and only if both $v$ and $-v$ are
regular values of $n$.
\end{lemma}

\begin{proof}\uses{lem:hess-height,lem:gauss-jacobian,lem:sign-kappa-index}
\eqref{eq:crit-split} is the first assertion of
Lemma~\ref{lem:hess-height}, the union being disjoint since $v\ne-v$.  By
Lemma~\ref{lem:gauss-jacobian}, $y$ is a regular value of $n$ iff
$\kappa\ne0$ on $n^{-1}(y)$, and by
Lemma~\ref{lem:sign-kappa-index} that is exactly nondegeneracy of the
corresponding critical points.
\end{proof}

\section{Morse functions and regular bands}\label{sec:morse}

\begin{lemma}[A Morse height function exists]\label{lem:morse-height}
\uses{cor:sard-charts,prop:orientation-parity,lem:critical-normal}
Let $\Sigma$ be a compact connected smooth embedded closed surface with the
unit normal $n$ supplied by
Proposition~\ref{prop:orientation-parity}\ref{it:parity-normal}.  Then for
almost every $v\in\Sph^2$ both $v$ and $-v$ are regular values of $n$, and for
every such $v$ the height function $h_v$ is Morse with finitely many critical
points, all lying in $n^{-1}(\{v,-v\})$.
\end{lemma}

\begin{proof}\uses{cor:sard-charts,lem:critical-normal}
Corollary~\ref{cor:sard-charts} applied to $n$, a $C^1$ map between
two-surfaces, makes the set of critical values $\Hs^2$-null; the union of that
set with its antipodal image is still null.  For $v$ outside it,
Lemma~\ref{lem:critical-normal} makes $h_v$ Morse, and its critical set is
then discrete and compact, hence finite.
\end{proof}

\begin{lemma}[Distinct critical values]\label{lem:morse-distinct}
\uses{lem:morse-height}
Let $h$ be Morse on the compact surface $\Sigma$.  There is a Morse function
$\widetilde h$ on $\Sigma$ with the same critical points, the same Hessians
there --- hence the same index counts $c_0,c_1,c_2$ --- and with pairwise
distinct critical values.
\end{lemma}

\begin{proof}\uses{lem:morse-height}
Write the critical points as $p_1,\dots,p_k$, choose disjoint neighbourhoods
$W_i\Subset V_i$ and smooth $b_i$ supported in $V_i$ with $b_i\equiv1$ on
$W_i$, and shrink the $W_i$ so that
$m:=\min_{\Sigma\setminus\bigcup_iW_i}|\nabla_\Sigma h|>0$.  Choose
arbitrarily small $c=(c_1,\dots,c_k)$ avoiding the finitely many affine
hyperplanes $h(p_i)+c_i=h(p_j)+c_j$ and small enough that
\begin{equation}\label{eq:morse-perturb}
  \sum_i|c_i|\,\|\nabla_\Sigma b_i\|_\infty<m/2 .
\end{equation}
Put $\widetilde h:=h+\sum_ic_ib_i$.  On each $W_i$ it differs from $h$ by the
additive constant $c_i$, so its critical points and Hessians there are those
of $h$; \eqref{eq:morse-perturb} prevents a critical point on
$\Sigma\setminus\bigcup_iW_i$, so none is created; and the critical values
$h(p_i)+c_i$ are distinct.
\end{proof}

\begin{fnote}[Why ``same Hessians'' is in the statement]\label{fnote:morse-hessian}
Proposition~\ref{prop:index-sum-antipodal} reads the index counts
$c_0,c_1,c_2$ off $h_v$ itself, while
Proposition~\ref{prop:morse-count} needs distinct critical values.  The bridge
is that Lemma~\ref{lem:morse-distinct} changes neither the critical points nor
their Hessians, hence neither the indices.  That is automatic from the
construction, but it has to be part of the statement rather than left inside
the proof.
\end{fnote}

\begin{lemma}[Local Morse coordinates]\label{lem:morse-coords}
Let $g$ be smooth near $0\in\R^2$ with $g(0)=0$, $Dg(0)=0$ and $D^2g(0)$
nonsingular of index $\lambda$.  There are smooth local coordinates in which
\begin{equation}\label{eq:morse-normal-form}
  g=\begin{cases}
    x_1^2+x_2^2,&\lambda=0,\\
    x_1^2-x_2^2,&\lambda=1,\\
    -x_1^2-x_2^2,&\lambda=2 .
  \end{cases}
\end{equation}
\end{lemma}

\begin{proof}
Taylor's integral formula writes
\begin{equation}\label{eq:morse-B}
  g(x)=\tfrac12x^{\mathsf T}B(x)x,
  \qquad
  B(x)=2\int_0^1(1-t)D^2g(tx)\,dt ,
\end{equation}
with $B(0)=D^2g(0)$.  After a fixed linear change a diagonal entry of $B(0)$
is nonzero; complete its square, use the resulting coordinate change (whose
Jacobian at zero is invertible), and repeat on the Schur complement.
\end{proof}

\begin{lemma}[Classification of a one-manifold with a tangent field]
\label{lem:one-manifold}
\uses{conv:ambient}
Let $M$ be a connected Hausdorff, second-countable smooth one-manifold without
boundary and suppose a smooth nowhere-zero tangent vector field on $M$ is
specified.  Then $M$ is
diffeomorphic to $S^1$ or to $\R$.  Consequently a noncompact such $M$ has
exactly two ends.  In particular this conclusion applies to every connected
component of a regular level set of a smooth function on an oriented
surface.
\end{lemma}

\begin{proof}\uses{thm:global-flow,cor:flow-manifold}
Normalize the specified field to unit length and use
Corollary~\ref{cor:flow-manifold} to take its maximal integral curve through a
base point.  Every orbit is open by the flow-box theorem, so the chosen orbit
and the union of all other orbits are disjoint open sets; connectedness makes
the chosen orbit all of $M$.  If its orbit map is injective, it is a bijective
local diffeomorphism from its maximal open interval onto $M$, hence a
diffeomorphism; every open interval is diffeomorphic to $\R$.  Otherwise the
curve returns to its base point.  ODE uniqueness then makes it periodic and
extends it for all time.  Its period group is closed and, by a flow box,
contains no nonzero sequence tending to zero; it is therefore $\tau\Z$ for a
least $\tau>0$, and the orbit map descends to a diffeomorphism
$\R/\tau\Z\simeq S^1\to M$.  For a regular level set on an oriented surface,
rotation by $\pi/2$ of the nonzero tangential gradient supplies the required
nowhere-zero tangent field.
\end{proof}

\begin{lemma}[Regular bands are smooth products]\label{lem:morse-band}
\uses{thm:global-flow,cor:flow-manifold}
Let $h$ be smooth on the compact surface $\Sigma$ and let $a<b$ be such that
$h^{-1}([a,b])$ contains no critical point.  Then
\begin{equation}\label{eq:band-product}
  \Theta:h^{-1}(a)\times[0,b-a]\to h^{-1}([a,b]),
  \qquad
  \Theta(p,t):=\Phi_t(p)
\end{equation}
is a diffeomorphism, where $\Phi$ is the flow of
\begin{equation}\label{eq:band-field}
  X:=\frac{\nabla_\Sigma h}{|\nabla_\Sigma h|^2}.
\end{equation}
\end{lemma}

\begin{proof}\uses{thm:global-flow,cor:flow-manifold,thm:flow-Ck}
$X$ is smooth on an open neighbourhood of the band, and
\begin{equation}\label{eq:band-derivative}
  \tfrac{d}{dt}h(\Phi_t(p))=dh(X)=1 .
\end{equation}
A trajectory starting on $h^{-1}(a)$ and run for $0\le t\le b-a$ remains in
the compact band, so the finite-step continuation of
Theorem~\ref{thm:global-flow} rules out a shorter maximal interval.  Hence
\eqref{eq:band-product} is a smooth bijection, with explicit smooth inverse
\begin{equation}\label{eq:band-inverse}
  q\longmapsto\bigl(\Phi_{a-h(q)}(q),\,h(q)-a\bigr).
\end{equation}
\end{proof}

\begin{corollary}[Sublevel retraction across a critical-value-free interval]
\label{cor:sublevel-stable}
\uses{lem:morse-band}
Let $h$ be Morse on the compact surface $\Sigma$ and let $a<b$.
\begin{enumerate}[label=(\roman*)]
\item\label{it:retract-closed} If $h$ has no critical value in $(a,b]$, then
  $\{h\le a\}$ is a deformation retract of $\{h\le b\}$.
\item\label{it:retract-open} If $h$ has no critical value in $[a,b)$, then
  $\{h\le a\}$ is a deformation retract of $\{h<b\}$.
\end{enumerate}
In either case the two sets have the same number of connected components.
\end{corollary}

\begin{proof}\uses{lem:morse-band,lem:morse-coords}
In case \ref{it:retract-closed}, $h^{-1}((a,b])$ contains no critical point, so
the flow $\Phi$ of \eqref{eq:band-field} is defined there and, by
\eqref{eq:band-derivative}, a trajectory started at $q$ with
$a<h(q)\le b$ exists backwards until its height tends to $a$.  If its limiting
point is regular, the ordinary local flow extends it to time $a-h(q)$.  Near a
critical point $p$ of level $a$, Morse coordinates and equivalence of the
coordinate and Riemannian norms give, on a sufficiently small coordinate ball,
\begin{equation}\label{eq:morse-flow-bounds}
 |\nabla_\Sigma h(x)|\ge c|x|,
 \qquad |h(x)-a|\le C|x|^2,
 \qquad
 |X(x)|\le \frac{C'}{\sqrt{h(x)-a}}\quad(h(x)>a).
\end{equation}
Along a descending trajectory, height itself is the time variable.  Thus a
segment descending from height $a+\delta$ has length at most
\begin{equation}\label{eq:morse-flow-length}
 \int_0^\delta \frac{C'}{\sqrt{\tau}}\,d\tau=2C'\sqrt\delta.
\end{equation}
A bootstrap with two concentric coordinate balls shows that a trajectory
starting sufficiently near $p$ cannot leave the larger ball before reaching
height $a$: otherwise its length would have a fixed positive lower bound,
contrary to \eqref{eq:morse-flow-length}.  The same estimate makes it Cauchy at
the finite terminal time and makes its endpoint tend to $p$ as $q\to p$.
Consequently every trajectory has a well-defined endpoint on level $a$, and
\begin{equation}\label{eq:retraction}
  r(q):=\begin{cases}q,&h(q)\le a,\\ \Phi_{a-h(q)}(q),&a<h(q)\le b\end{cases}
\end{equation}
is a well-defined retraction of $\{h\le b\}$ onto $\{h\le a\}$.  It is
continuous at regular points of $h^{-1}(a)$ because $\Phi_0=\Id$, and
\eqref{eq:morse-flow-length} proves continuity at a critical point of that
level.  The identical length estimate, applied to partial trajectory segments,
shows jointly in $(q,s)$ that
\[
 H_s(q):=\begin{cases}
 q,&h(q)\le a,\\
 \Phi_{s(a-h(q))}(q),&a<h(q)\le b,
 \end{cases}
 \qquad 0\le s\le1,
\]
is continuous and stays in $\{h\le b\}$; it is a deformation from the
identity to $r$.  Thus every component meets $\{h\le a\}$ and two such
components cannot be joined upstairs without being joined after applying
$r$, so the component counts agree.  Case \ref{it:retract-open} uses the
same formulas with $h(q)<b$.
\end{proof}

\begin{lemma}[Local structure at a critical point]\label{lem:local-sectors}
\uses{lem:morse-coords}
Let $h$ be Morse on the surface $\Sigma$, let $p$ be a critical point of index
$\lambda$ and $c:=h(p)$, and let $D$ be a Morse coordinate ball about $p$ as
in Lemma~\ref{lem:morse-coords}.  Put $L:=h^{-1}(c)$.
\begin{enumerate}[label=(\roman*)]
\item\label{it:sect-0} $\lambda=0$: $\{h<c\}\cap D=\varnothing$ and
  $\{h>c\}\cap D=D\setminus\{p\}$, and $L\cap D=\{p\}$.
\item\label{it:sect-2} $\lambda=2$: $\{h>c\}\cap D=\varnothing$ and
  $\{h<c\}\cap D=D\setminus\{p\}$, and $L\cap D=\{p\}$.
\item\label{it:sect-1} $\lambda=1$: $L\cap D$ is the union of four smooth arcs
  (\emph{rays}) $r_1,r_2,r_3,r_4$ from $p$ to $\partial D$, meeting only at
  $p$; $\{h<c\}\cap D$ and $\{h>c\}\cap D$ each have exactly two connected
  components (\emph{sectors}); and in cyclic order around $p$ the rays and
  sectors alternate as
  \[
    r_1,\ S^-_{12},\ r_2,\ S^+_{23},\ r_3,\ S^-_{34},\ r_4,\ S^+_{41},
  \]
  where $S^-_{\bullet}\subset\{h<c\}$ and $S^+_{\bullet}\subset\{h>c\}$.  In
  particular the two sectors flanking $r_1$ are $S^+_{41}$ and $S^-_{12}$,
  those flanking $r_2$ are $S^-_{12}$ and $S^+_{23}$, those flanking $r_3$ are
  $S^+_{23}$ and $S^-_{34}$, and those flanking $r_4$ are $S^-_{34}$ and
  $S^+_{41}$.
\end{enumerate}
\end{lemma}

\begin{proof}\uses{lem:morse-coords}
Read \eqref{eq:morse-normal-form} off in the chart.  For $\lambda=1$ the set
$\{x_1^2=x_2^2\}$ is the union of the four rays
$x_2=\pm x_1$, $\pm x_1>0$; $\{x_1^2<x_2^2\}$ is the pair of vertical sectors
and $\{x_1^2>x_2^2\}$ the pair of horizontal ones; and the cyclic alternation
is immediate.
\end{proof}

\begin{fnote}[No handle attachment is used]\label{fnote:no-handles}
Lemma~\ref{lem:local-sectors} is purely local: it says nothing relating
$h^{-1}(c-\varepsilon)$ to $h^{-1}(c+\varepsilon)$.  An earlier draft of this
section asserted such a relation --- that the level set is obtained by deleting
two arcs and inserting two others --- and used it to run a surgery argument.
That assertion is true but its proof is the handle-attachment theorem, which
requires constructing a modified function with a bump near $p$ and is a
substantial development in its own right.  \S\ref{sec:morse-count} is written
so that it is never needed: every comparison there is between $\{h<c\}$,
$\{h\le c\}$ and a sublevel set at a nearby regular value, all handled by
Lemma~\ref{lem:local-sectors} and
Corollary~\ref{cor:sublevel-stable}.
\end{fnote}

\section{The Morse count is at most two}\label{sec:morse-count}

\begin{definition}[Merging saddle]\label{def:merging-saddle}
\uses{lem:local-sectors}
Let $h$ be Morse on the compact connected surface $\Sigma$ and let $p$ be a
critical point of index $1$ with $c:=h(p)$.  With the notation of
Lemma~\ref{lem:local-sectors}\ref{it:sect-1}, write $A(S)$ for the connected
component of $\{h<c\}$ containing a sector $S$, and $B(S)$ for the component
of $\{h>c\}$ containing $S$.  Say $p$ is
\begin{itemize}[nosep]
\item \emph{sub-merging} if $A(S^-_{12})\ne A(S^-_{34})$;
\item \emph{super-merging} if $B(S^+_{23})\ne B(S^+_{41})$.
\end{itemize}
\end{definition}

\begin{lemma}[Adjacent components are constant along a regular level arc]
\label{lem:level-adjacency}
\uses{conv:ambient}
Let $h$ be smooth on $\Sigma$, $c\in\R$, $L:=h^{-1}(c)$, and let $q\in L$ be a
regular point of $h$.  Then $q$ has a neighbourhood $U$ with
$U\cap\{h<c\}$ and $U\cap\{h>c\}$ both connected and nonempty; writing
$\phi^-(q)$ and $\phi^+(q)$ for the components of $\{h<c\}$ and $\{h>c\}$
containing them, the maps $\phi^\pm$ are locally constant on the set of
regular points of $L$, hence constant on each connected subset of it.
\end{lemma}

\begin{proof}
By the implicit function theorem \ref{base:calculus}, in suitable coordinates
$h-c$ is the first coordinate on a ball $U$; then $U\cap\{h<c\}$ and
$U\cap\{h>c\}$ are the two open half-balls, each connected and nonempty.  For
$q'\in U\cap L$ the same half-balls witness
$\phi^\pm(q')=\phi^\pm(q)$.
\end{proof}

\begin{lemma}[Closing up a sublevel set]\label{lem:sublevel-closure}
\uses{def:merging-saddle,lem:local-sectors,lem:level-adjacency}
Let $h$ be Morse with distinct critical values on the compact connected
surface $\Sigma$, let $c$ be a critical value with unique critical point $p$,
and let $\lambda$ be its index.  Then
\[
  \#\mathrm{comp}\{h\le c\}-\#\mathrm{comp}\{h<c\}
  =\begin{cases}
    +1,&\lambda=0,\\
    -1,&\lambda=1\text{ and }p\text{ sub-merging},\\
    0,&\text{otherwise}.
  \end{cases}
\]
\end{lemma}

\begin{proof}\uses{lem:local-sectors,lem:level-adjacency,def:merging-saddle}
Write $L:=h^{-1}(c)$, so $\{h\le c\}=\{h<c\}\sqcup L$.  A component of
$\{h\le c\}$ is a union of components of $\{h<c\}$ together with part of $L$,
so it suffices to determine which components of $\{h<c\}$ get identified and
whether any part of $L$ forms a component by itself.

Away from a smaller Morse coordinate ball around $p$, the compact regular
level set is a compact one-manifold with boundary and therefore has finitely
many components.  Inside the ball it consists of the pieces listed in
Lemma~\ref{lem:local-sectors}.  Thus $L$ is a finite graph whose only possible
nonmanifold vertex is $p$.  By Lemma~\ref{lem:level-adjacency}, $\phi^-$ is
constant on each edge in the regular part $L\setminus\{p\}$.  Consequently the
closure of an edge attaches along all of its regular points to one and the
same component of $\{h<c\}$; distinct such components can be identified by
adjoining $L$ only at the vertex $p$.  A closed regular-level component not
meeting $p$ likewise attaches to exactly one sublevel component and identifies
nothing.

At $p$: if $\lambda=0$ then by
Lemma~\ref{lem:local-sectors}\ref{it:sect-0} a neighbourhood of $p$ meets
$\{h<c\}$ not at all and meets $L$ only in $\{p\}$, so $\{p\}$ is an isolated
point of $\{h\le c\}$ and contributes one new component.  If $\lambda=2$ then
by \ref{it:sect-2} the punctured ball lies in $\{h<c\}$, so $p$ is adjacent to
exactly one component and nothing changes.  If $\lambda=1$ then by
\ref{it:sect-1} $p$ is adjacent to exactly the two sectors $S^-_{12}$ and
$S^-_{34}$, so it identifies $A(S^-_{12})$ with $A(S^-_{34})$: the count drops
by one precisely when these differ, i.e.\ precisely when $p$ is sub-merging.
No component of $\{h<c\}$ is created or destroyed otherwise.
\end{proof}

\begin{lemma}[Counting merging saddles]\label{lem:count-submerging}
\uses{lem:sublevel-closure,cor:sublevel-stable,def:merging-saddle}
Let $h$ be Morse with distinct critical values on the compact connected
surface $\Sigma$, with $c_0$ minima and $c_2$ maxima.  Then the number of
sub-merging saddles is exactly $c_0-1$, and the number of super-merging
saddles is exactly $c_2-1$.
\end{lemma}

\begin{proof}\uses{lem:sublevel-closure,cor:sublevel-stable}
List the critical values $c^{(1)}<\cdots<c^{(k)}$ and pick regular values
$a_0<c^{(1)}<a_1<c^{(2)}<\cdots<c^{(k)}<a_k$.  Put
$N_j:=\#\mathrm{comp}\{h\le a_j\}$.  Then $N_0=0$ since
$\{h\le a_0\}=\varnothing$, and $N_k=1$ since $\{h\le a_k\}=\Sigma$ is
connected.

Fix $j$ and write $c:=c^{(j)}$.  Corollary~\ref{cor:sublevel-stable}\ref{it:retract-open}
applied on $[a_{j-1},c)$ gives
$\#\mathrm{comp}\{h<c\}=N_{j-1}$, and
Corollary~\ref{cor:sublevel-stable}\ref{it:retract-closed} applied on
$(c,a_j]$ gives $N_j=\#\mathrm{comp}\{h\le c\}$.  So
Lemma~\ref{lem:sublevel-closure} computes $N_j-N_{j-1}$ as $+1$ at a minimum,
$-1$ at a sub-merging saddle, and $0$ otherwise.

Summing over $j$, $1-0=c_0-\#\{\text{sub-merging}\}$.  Apply the same argument
to $-h$: its minima are the maxima of $h$, its index-$1$ points are those of
$h$ because $\mathrm{ind}_p(-h)=2-\mathrm{ind}_p(h)$, and sub-merging for $-h$
is super-merging for $h$ because $\{-h<-c\}=\{h>c\}$.
\end{proof}

\begin{lemma}[No saddle merges on both sides]\label{lem:merge-disjoint}
\uses{def:merging-saddle,lem:local-sectors,lem:level-adjacency,lem:one-manifold}
Let $h$ be Morse with pairwise distinct critical values on the compact
connected surface $\Sigma$.  Then no saddle of $h$ is both sub-merging and
super-merging.
\end{lemma}

\begin{proof}\uses{lem:local-sectors,lem:level-adjacency,lem:one-manifold,
                    prop:orientation-parity}
Let $p$ be a saddle, $c:=h(p)$, $L:=h^{-1}(c)$, and adopt the notation of
Lemma~\ref{lem:local-sectors}\ref{it:sect-1}.  The distinct-critical-value
hypothesis makes $p$ the only critical point at level $c$; hence
$L\setminus\{p\}$ is a smooth one-manifold without boundary, by the
implicit function theorem \ref{base:calculus}.

Let $\sigma$ be the connected component of $L\setminus\{p\}$ containing
$r_1\setminus\{p\}$.  It is not compact, since $r_1$ accumulates at $p\notin\sigma$,
and it carries the nowhere-zero tangent field obtained by rotating
$\nabla_\Sigma h$ by $\pi/2$ using the orientation of $\Sigma$.  Hence
Lemma~\ref{lem:one-manifold} makes it diffeomorphic to $\R$ with exactly two
ends.  Both ends accumulate only at $p$: $L$ is compact and
$L\setminus\{p\}$ is a manifold, so any limit point of $\sigma$ outside
$\sigma$ lies in $L$ and is not a manifold point, hence equals $p$.  Near $p$
the set $L$ is the union of the four rays
(Lemma~\ref{lem:local-sectors}\ref{it:sect-1}), and each ray accounts for
exactly one end; the two ends of $\sigma$ therefore lie along two
\emph{distinct} rays, one of which is $r_1$.  Write $r_j$ for the other,
$j\in\{2,3,4\}$.

By Lemma~\ref{lem:level-adjacency}, $\phi^-$ and $\phi^+$ are constant on the
connected set $\sigma$.  Along the $r_1$ end, the flanking sectors are
$S^+_{41}$ and $S^-_{12}$, so
$\phi^-(\sigma)=A(S^-_{12})$ and $\phi^+(\sigma)=B(S^+_{41})$.  Reading the
flanking sectors of $r_j$ off Lemma~\ref{lem:local-sectors}\ref{it:sect-1}:
\begin{itemize}[nosep]
\item $j=2$: flanking sectors $S^-_{12}$, $S^+_{23}$, so
  $B(S^+_{41})=\phi^+(\sigma)=B(S^+_{23})$ and $p$ is not super-merging.
\item $j=4$: flanking sectors $S^-_{34}$, $S^+_{41}$, so
  $A(S^-_{12})=\phi^-(\sigma)=A(S^-_{34})$ and $p$ is not sub-merging.
\item $j=3$: flanking sectors $S^+_{23}$, $S^-_{34}$, so both equalities hold
  and $p$ is neither.
\end{itemize}
In every case at least one of the two merging properties fails.
\end{proof}

\medskip\noindent\emph{Formalisation note.} The formal proof follows the integral curve of $n\times\nabla_\Sigma h$ through a point of $r_1$ directly: it is injective, stays in $L\setminus\{p\}$, and tends to $p$ at both ends along two distinct rays, which gives the two-ends statement used above without invoking Lemma~\ref{lem:one-manifold}. That lemma is formalised independently, in its abstract form (second countability is not needed) and in the embedded form for level components.

\begin{proposition}[The Morse count is at most two]\label{prop:morse-count}
\uses{lem:count-submerging,lem:merge-disjoint}
Let $h$ be Morse on a compact connected surface $\Sigma$, with $c_\lambda$
critical points of index $\lambda$.  Then
\begin{equation}\label{eq:morse-count}
  c_0-c_1+c_2\le2 .
\end{equation}
\end{proposition}

\begin{proof}\uses{lem:count-submerging,lem:merge-disjoint,lem:morse-distinct}
By Lemma~\ref{lem:morse-distinct} we may assume the critical values are
distinct, since the perturbation changes no $c_\lambda$.  Then
Lemma~\ref{lem:count-submerging} produces $c_0-1$ sub-merging and $c_2-1$
super-merging saddles, and Lemma~\ref{lem:merge-disjoint} makes these two sets
disjoint, so $c_1\ge(c_0-1)+(c_2-1)$.
\end{proof}

\begin{remark}[What is \emph{not} claimed]\label{rem:no-euler}
Proposition~\ref{prop:morse-count} is an inequality, and that is all this
development needs.  No Euler characteristic is defined, no homology is
computed, no cell structure is built, no handle is attached, and the
classification of compact surfaces is not used.  In particular the identity
$c_0-c_1+c_2=2-2g$ is deliberately \emph{not} proved, and neither is
independence of $h$.
\end{remark}

\section{The total curvature bound}\label{sec:total-curvature}

\begin{lemma}[Regular fibres are finite]\label{lem:regular-fiber-finite}
\uses{conv:ambient}
Let $f:\Sigma\to\Sigma'$ be a $C^1$ map between compact surfaces, $\Sigma$
without boundary, and let $y$ be a regular value.  Then $f^{-1}(y)$ is finite.
\end{lemma}

\begin{proof}
The inverse function theorem \ref{base:calculus} makes $f^{-1}(y)$ discrete;
it is closed in the compact $\Sigma$, hence finite.
\end{proof}

\begin{definition}[Local index sum]\label{def:index-sum}
\uses{lem:regular-fiber-finite,prop:orientation-parity}
Let $\Sigma,\Sigma'$ be smooth surfaces with $\Sigma$ compact, without
boundary and oriented, $\Sigma'$ oriented, and let $f:\Sigma\to\Sigma'$ be
$C^1$.  For a regular value $y\in\Sigma'$ put
$\varepsilon_f(p):=\operatorname{sign}\det d_pf$ using the orientations, and
\begin{equation}\label{eq:index-sum}
  \iota_y(f):=\sum_{p\in f^{-1}(y)}\varepsilon_f(p)\ \in\Z ,
\end{equation}
a finite sum by Lemma~\ref{lem:regular-fiber-finite}.  At a nonregular value
define $\iota_y(f):=0$.  Thus $y\mapsto\iota_y(f)$ is an everywhere-defined
integer-valued function.
\end{definition}

\begin{proof}\uses{lem:regular-fiber-finite}
The critical set is closed in the compact source, so its image is compact;
hence the regular-value set is open.  At a regular value the fibre is finite.
Choose pairwise disjoint inverse-function neighbourhoods of its points; by
compactness, after shrinking the common target neighbourhood, its full
preimage lies in their union.  The signs of the local Jacobians are constant
there, so $\iota_y(f)$ is locally constant on the regular-value set.  Its
extension by zero over the closed critical-value set is therefore Borel.
\end{proof}

\begin{fnote}[This is not a degree]\label{fnote:not-a-degree}
$\iota_y(f)$ is \emph{not} asserted to be independent of $y$, and no such
statement is proved or used.  Proving independence would require a
transversality statement for a path in $\Sigma'$ against $f$, which
Theorem~\ref{thm:transversality} does \emph{not} supply --- it is about maps
into $\R^3$ transverse to a surface in $\R^3$.  The antipodal averaging in
Theorem~\ref{thm:total-curvature-bound} is what makes independence
unnecessary; see Formalisation note~\ref{fnote:why-averaging}.
\end{fnote}

\begin{lemma}[The Gauss map extends]\label{lem:gauss-map-extension}
\uses{thm:tubular}
Let $\Sigma\subset\R^3$ be a compact smooth embedded closed surface with unit
normal $n$.  There is a Lipschitz $N:\R^3\to\R^3$ with $N|_\Sigma=n$ and
$D N(p)|_{T_p\Sigma}=d_pn$ for every $p\in\Sigma$.
\end{lemma}

\begin{proof}\uses{thm:tubular}
Theorem~\ref{thm:tubular} gives a tubular neighbourhood
$\Theta(p,t)=p+tn(p)$ and hence a smooth nearest-point projection
$\pi$ onto $\Sigma$ there.  Put $N:=\chi\cdot(n\circ\pi)$ with $\chi$ a smooth
cutoff equal to one near $\Sigma$ and supported in the tubular neighbourhood.
Then $N|_\Sigma=n$, and since $\pi|_\Sigma=\mathrm{id}$ and
$d_p\pi|_{T_p\Sigma}=\mathrm{id}$, the tangential differential is $d_pn$.
\end{proof}

\begin{theorem}[Total curvature as an integral of index sums]
\label{thm:total-curvature-index}
\uses{lem:gauss-jacobian,def:index-sum,thm:area-formula-rect,
      prop:orientation-parity}
Let $\Sigma\subset\R^3$ be a compact connected smooth embedded closed surface
with the unit normal $n$ of
Proposition~\ref{prop:orientation-parity}\ref{it:parity-normal}.  Then
\begin{equation}\label{eq:total-curvature-index}
  \int_\Sigma\kappa\,d\Hs^2=\int_{\Sph^2}\iota_y(n)\,d\Hs^2(y),
\end{equation}
both sides being finite.
\end{theorem}

\begin{proof}
\uses{lem:gauss-jacobian,def:index-sum,thm:area-formula-rect,
      cor:sard-charts,lem:gauss-map-extension,
      lem:smooth-surface-rectifiable,conv:gauss-orientation,lem:sphere-area}
By Lemma~\ref{lem:smooth-surface-rectifiable} and
Lemma~\ref{lem:gauss-map-extension} the area formula
\eqref{eq:area-formula-rect} applies to $n$ on the rectifiable set $\Sigma$,
with $J_2(d_pn)=|\kappa(p)|$ by \eqref{eq:gauss-jacobian}.  Apply it twice,
with the nonnegative Borel weights
\[
  q_\pm:=\max\{\pm\operatorname{sign}\kappa,0\},
\]
and subtract; both integrals are finite because $\kappa$ is continuous on the
compact $\Sigma$.  Since $q_\pm J_2(dn)=\kappa^\pm$, the left side becomes
$\int_\Sigma\kappa$, and the right side becomes
\[
  \int_{\Sph^2}\Bigl(\sum_{p\in n^{-1}(y),\ \kappa(p)>0}1
    -\sum_{p\in n^{-1}(y),\ \kappa(p)<0}1\Bigr)d\Hs^2(y).
\]
By Corollary~\ref{cor:sard-charts} almost every $y$ is a regular value of $n$,
and at a regular value every $p\in n^{-1}(y)$ has $\kappa(p)\ne0$ with
$\varepsilon_n(p)=\operatorname{sign}\kappa(p)$ by
\eqref{eq:gauss-orientation}; so for almost every $y$ the bracket is
$\iota_y(n)$.  The right-hand integral is supported on $\Sph^2$, whose total
$\Hs^2$-measure is $4\pi$ by Lemma~\ref{lem:sphere-area}.  At the exceptional
critical values no pointwise identification with $\iota_y(n)$ is needed,
because that target set is $\Hs^2$-null.
\end{proof}

\medskip\noindent\emph{Formalisation note.} The identity is formalised for every compact smooth embedded surface and every unit normal field; connectedness and the particular normal of Proposition~\ref{prop:orientation-parity} are not needed.

\begin{proposition}[Antipodal index sums add up to the Morse count]
\label{prop:index-sum-antipodal}
\uses{def:index-sum,lem:sign-kappa-index,lem:critical-normal,lem:morse-height,
      prop:orientation-parity}
Let $\Sigma\subset\R^3$ be a compact connected smooth embedded closed
surface, let $n$ be the unit normal supplied by
Proposition~\ref{prop:orientation-parity}, and let $v\in\Sph^2$ be such that
both $v$ and $-v$ are regular values of $n$ --- almost every $v$, by
Lemma~\ref{lem:morse-height}.  Let $c_\lambda$ be the number of critical
points of $h_v$ of index $\lambda$.  Then
\begin{equation}\label{eq:index-sum-antipodal}
  \iota_v(n)+\iota_{-v}(n)=c_0-c_1+c_2 .
\end{equation}
\end{proposition}

\begin{proof}
\uses{def:index-sum,lem:sign-kappa-index,lem:critical-normal,
      conv:gauss-orientation}
By \eqref{eq:gauss-orientation},
$\varepsilon_n(p)=\operatorname{sign}\kappa(p)$, and by
Lemma~\ref{lem:sign-kappa-index} this equals $(-1)^{\mathrm{ind}_p(h_v)}$ at
every critical point of $h_v$.  Hence
\[
  \iota_v(n)=\!\!\sum_{p\in n^{-1}(v)}\!\!(-1)^{\mathrm{ind}_p(h_v)},
  \qquad
  \iota_{-v}(n)=\!\!\sum_{p\in n^{-1}(-v)}\!\!(-1)^{\mathrm{ind}_p(h_v)} .
\]
Adding these and using the disjoint splitting \eqref{eq:crit-split} of
Lemma~\ref{lem:critical-normal} gives \eqref{eq:index-sum-antipodal}.
\end{proof}

\begin{theorem}[Total curvature bound]\label{thm:total-curvature-bound}
\uses{thm:total-curvature-index,prop:index-sum-antipodal,prop:morse-count}
Let $\Sigma\subset\R^3$ be a compact connected smooth embedded closed
surface.  Then
\begin{equation}\label{eq:total-curvature-bound}
  \int_\Sigma\kappa\,d\Hs^2\le4\pi .
\end{equation}
\end{theorem}

\begin{proof}
\uses{thm:total-curvature-index,prop:index-sum-antipodal,prop:morse-count,
      lem:morse-height,prop:orientation-parity,lem:sphere-area}
Orient $\Sigma$ and fix $n$ by
Proposition~\ref{prop:orientation-parity}\ref{it:parity-normal}.  The
antipodal map is an isometry of $\Sph^2$, so it preserves $\Hs^2$ and
\[
  \int_{\Sph^2}\iota_y(n)\,d\Hs^2(y)
  =\frac12\int_{\Sph^2}\bigl(\iota_y(n)+\iota_{-y}(n)\bigr)\,d\Hs^2(y).
\]
For almost every $y$ both $y$ and $-y$ are regular values of $n$
(Lemma~\ref{lem:morse-height}), and for such $y$
Proposition~\ref{prop:index-sum-antipodal} and
Proposition~\ref{prop:morse-count} applied to the Morse function $h_y$ give
\[
  \iota_y(n)+\iota_{-y}(n)=c_0-c_1+c_2\le2 .
\]
Hence, with \eqref{eq:total-curvature-index} and
$\Hs^2(\Sph^2)=4\pi$,
\[
  \int_\Sigma\kappa\,d\Hs^2
  =\int_{\Sph^2}\iota_y(n)\,d\Hs^2(y)
  \le\tfrac12\cdot2\cdot4\pi=4\pi .
\]
\end{proof}

\begin{fnote}[Why averaging replaces degree theory]\label{fnote:why-averaging}
The classical route asserts $\iota_y(n)=\deg(n)$ for every regular $y$ and
then bounds $\deg(n)\le1$.  Establishing $y$-independence costs a
transversality statement this document does not have
(Formalisation note~\ref{fnote:not-a-degree}).  It is also unnecessary: the
quantity actually needed is the \emph{average} of $\iota_y(n)$ over
$\Sph^2$, and the pointwise bound
$\iota_y(n)+\iota_{-y}(n)\le2$ from
Proposition~\ref{prop:index-sum-antipodal} bounds that average by one
directly.  This removes an entire section from the chapter and, with it, the
only place where orientations of one-manifold cobordisms would have had to be
tracked.
\end{fnote}

\begin{fnote}[Why this route]\label{fnote:why-degree}
An earlier design proved \eqref{eq:total-curvature-bound} from a finite smooth
triangulation, a local Gauss--Bonnet formula on curved triangles, and
$\chi=b_0-b_1+b_2\le2$ via the cellular chain complex.  That route required,
in addition to everything above: abstract simplicial complexes with geometric
realisation, a four-part simplicial gluing lemma with a quotient-complex
argument, a moving-frame connection and structure equation, Stokes on a curved
triangle with corners, the turning-number calculation (polar decomposition and
angle lifting over the disk), and the $b_2=1$ dual-graph argument.  The pinned
mathlib supplies none of it: \code{AlgebraicTopology/Simplicial*} is
category-theoretic simplicial sets and
\code{Analysis/Convex/SimplicialComplex} is convex geometry, so all of it
would have been built from zero.

What replaces it is \S\ref{sec:shape} (pointwise linear algebra),
\S\ref{sec:morse-count} (component counting plus
Lemma~\ref{lem:merge-disjoint}) and \S\ref{sec:total-curvature} (the area
formula plus antipodal averaging).  The exchange is not free: the sharp
\emph{identity} $\int\kappa=2\pi\chi$ is lost, and only the inequality
\eqref{eq:total-curvature-bound} survives
(Remark~\ref{rem:no-euler}).  Since
Lemma~\ref{lem:K-gauss-bonnet-input} is the sole consumer and needs only the
inequality, that costs nothing here.  It would matter if this development were
reused as a general-purpose surface-theory library, which it is not intended
to be.
\end{fnote}

\begin{fnote}[The remaining cost centre]\label{fnote:retire-flows}
After the exchange described in Formalisation note~\ref{fnote:why-degree}, the
dominant cost of this chapter is \S\ref{sec:flow} together with
\S\ref{sec:morse}: the global flow is needed only for
Lemma~\ref{lem:morse-band} and Lemma~\ref{lem:one-manifold}, which feed only
Corollary~\ref{cor:sublevel-stable} and Lemma~\ref{lem:merge-disjoint}, which
in turn feed \S\ref{sec:morse-count}.  The shape-operator calculations and
the index-integral formula can be proved independently of that flow
construction.  The final total-curvature bound uses
Proposition~\ref{prop:morse-count}, so the complete
\code{Surface/TotalCurvature.lean} file depends transitively on
\texttt{Flow/}.  Two observations for anyone trying to reduce that dependency
further.
\begin{enumerate}[label=(\alph*)]
\item Corollary~\ref{cor:sublevel-stable} is used only for the \emph{count} of
  connected components, never for the retraction itself.  Its
  \ref{it:retract-open} half already has a flow-free proof of surjectivity ---
  a component $V$ of $\{h<b\}$ has $\min_{\overline V}h$ attained at an
  interior point, hence at a critical point of value $<b$ --- and only
  injectivity uses the flow.  If injectivity can be obtained without it,
  \S\ref{sec:flow} would be needed only for
  Lemma~\ref{lem:one-manifold}, whose sole use is to know that a connected
  one-manifold has two ends.  I have verified neither step.
\item On a level set $\Sigma_t=\{u=t\}$ of the capacitary potential the Gauss
  map \emph{is} $-\nabla u/|\nabla u|$, which extends to the whole exterior
  region, and by \eqref{eq:grad-u-expansion} it has degree one on a large
  sphere.  A degree comparison between $\Sigma_t$ and that sphere would give
  $\deg(n)\le1$ directly, retiring \S\ref{sec:morse} \emph{and}
  \S\ref{sec:flow}.  The obstruction is that the critical set of a harmonic
  function need not be discrete --- precisely the difficulty
  Chapter~\ref{ch:appK} is designed to route around --- so the comparison
  would have to be run with that set excised.  Whether that can be done is
  open; do not plan on it.
\end{enumerate}
\S\ref{sec:transversality}, \S\ref{sec:tubular} and
Theorem~\ref{thm:component-count} cannot be retired under any of these
variants: Lemma~\ref{lem:hull-properties} and
Lemma~\ref{lem:level-connected} use them independently of this chapter.
\end{fnote}

